% ============================================================
% preamble.tex —— 《SYZ 猜想综述集》导言区（Lead 独占；译者只读）
% 编译：LuaLaTeX ×3（ctexbook，A4，11pt）
% ============================================================

% ---------- 数学 ----------
\usepackage{amsmath,amssymb,amsthm,mathtools}
\usepackage{bm}
\usepackage{mathrsfs}

% ---------- 排版 ----------
\usepackage[a4paper,top=2.40cm,bottom=2.40cm,left=2.60cm,right=2.60cm,headheight=15pt]{geometry}
\usepackage{graphicx}
\usepackage{xcolor}
\usepackage{enumitem}
\usepackage{array,booktabs,multirow,longtable,tabularx}
\usepackage{multicol}
\usepackage{fancyhdr}
\usepackage{tcolorbox}
\usepackage{tikz}
\usetikzlibrary{arrows.meta,calc,cd,decorations.pathmorphing,decorations.markings,positioning,shapes.geometric,shapes.misc,shapes.arrows,patterns,intersections,through,backgrounds}
\usepackage{caption}
\usepackage[hidelinks,bookmarksnumbered]{hyperref}

\allowdisplaybreaks
\linespread{1.18}
\setlength{\parskip}{0.25em}
\setlength{\emergencystretch}{3em}% 降低 Overfull hbox（中文长行/公式混合）
\setlist[itemize]{leftmargin=2.2em,topsep=0.35em,itemsep=0.15em,parsep=0.15em}
\setlist[enumerate]{leftmargin=2.6em,topsep=0.35em,itemsep=0.15em,parsep=0.15em}

% ---------- 中文标题样式 ----------
\ctexset{
  chapter/name   = {第,章},
  chapter/format = \Large\bfseries,
  section/format = \large\bfseries,
  subsection/format = \normalsize\bfseries,
  subsubsection/format = \normalsize\bfseries,
}

% ---------- 页眉页脚 ----------
\pagestyle{fancy}
\fancyhf{}
\fancyhead[L]{\footnotesize\leftmark}
\fancyhead[R]{\footnotesize\rightmark}
\fancyfoot[C]{\small\thepage}
\renewcommand{\headrulewidth}{0.4pt}
\renewcommand{\chaptermark}[1]{\markboth{第\thechapter\ 章\quad #1}{}}
\renewcommand{\sectionmark}[1]{\markright{\thesection\quad #1}}
\fancypagestyle{plain}{%
  \fancyhf{}%
  \fancyfoot[C]{\small\thepage}%
  \renewcommand{\headrulewidth}{0pt}%
}

% ---------- 论文出处块（每章开头，Lead 写）----------
\newcommand{\papersource}[4]{%
  \par\medskip
  \begin{flushleft}\small
    {\bfseries 英文原题：}#1\\[2pt]
    {\bfseries 作者：}#2\\[2pt]
    {\bfseries 出处：}#3\\[2pt]
    {\bfseries 底本页数：}#4
  \end{flushleft}
  \medskip
}

% ---------- 定理类环境（编号手工给定，如 \begin{theorem}{4.5.2}）----------
% 可选第二参数为命名：\begin{theorem}{4.5.2}[Hein]
\newcommand{\DeclareThmEnv}[2]{%
  \NewDocumentEnvironment{#1}{m o}
    {\par\medskip\noindent\textbf{#2 ##1\IfValueT{##2}{（##2）}}\hspace{0.5em}}
    {\par\medskip}%
}
\DeclareThmEnv{theorem}{定理}
\DeclareThmEnv{proposition}{命题}
\DeclareThmEnv{lemma}{引理}
\DeclareThmEnv{corollary}{推论}
\DeclareThmEnv{definition}{定义}
\DeclareThmEnv{conjecture}{猜想}
\DeclareThmEnv{question}{问题}
\DeclareThmEnv{remark}{注}
\DeclareThmEnv{example}{例}
\DeclareThmEnv{observation}{观察}
\DeclareThmEnv{claim}{断言}
\DeclareThmEnv{hypothesis}{假设}
\DeclareThmEnv{construction}{构造}
\DeclareThmEnv{fact}{事实}
\DeclareThmEnv{notation}{记号}
\DeclareThmEnv{setup}{约定}
\DeclareThmEnv{exercise}{习题}
\DeclareThmEnv{answer}{解答}

\let\proof\relax
\let\endproof\relax
\newenvironment{proof}[1][证明]{\par\medskip\noindent\textit{#1}\ \ignorespaces}{\par\medskip}
\newcommand{\proofend}{\hfill$\square$}

% ---------- 行内小标题（原文 \textbf{Tame punctures.} 之类）----------
\newcommand{\para}[1]{\par\medskip\noindent\textbf{#1}\ \ignorespaces}

% ---------- 图 / 表 题注（手工编号；\figimg 自闭合，\figtitle 紧随其后）----------
\newcommand{\figimg}[2][0.80]{%
  \begin{center}\includegraphics[width=#1\textwidth]{#2}\end{center}\nopagebreak[4]}
\newcommand{\figtitle}[2]{%
  \par\nopagebreak[4]\vspace{2pt}\begin{center}\small\textbf{图 #1.}\quad #2\end{center}\par\smallskip}
\newcommand{\tabtitle}[2]{%
  \par\smallskip
  \begin{center}\small\textbf{表 #1.}\quad #2\end{center}
  \par\smallskip}

% ---------- 要点框 ----------
\newtcolorbox{keybox}[1]{%
  colback=gray!6, colframe=gray!55, coltitle=black, colbacktitle=gray!20,
  title={#1}, fonttitle=\bfseries, boxrule=0.4pt, arc=2pt,
  left=6pt, right=6pt, top=4pt, bottom=4pt, before skip=6pt, after skip=6pt}
\newtcolorbox{keyboxplain}{%
  colback=gray!6, frame hidden, boxrule=0pt, arc=2pt,
  left=6pt, right=6pt, top=4pt, bottom=4pt, before skip=6pt, after skip=6pt}

% ---------- 参考文献条目（gen_refs.py 生成；英文原貌）----------
\newcommand{\refitem}[2]{\par\hangindent=1.6em\hangafter=1\textbf{[#1]}\hspace{0.45em}#2\par}

% ---------- 常用符号（按需增补；译者可直接用 \mathbb 等）----------
\newcommand{\R}{\mathbb{R}}
\newcommand{\C}{\mathbb{C}}
\newcommand{\Z}{\mathbb{Z}}
\newcommand{\Q}{\mathbb{Q}}
\newcommand{\N}{\mathbb{N}}
\newcommand{\PP}{\mathbb{P}}
\newcommand{\ii}{\mathrm{i}}
\newcommand{\ee}{\mathrm{e}}
\newcommand{\dd}{\mathrm{d}}
\newcommand{\Log}{\operatorname{Log}}
\graphicspath{{figures/}}
